% Môn: Toán
% Lớp: 10
% Chương: Chương I. Mệnh đề và tập hợp
% Bài: Bài 2. Tập hợp và các phép toán trên tập hợp
% Số câu: 56
% App đọc file này trực tiếp — sửa rồi Commit trên GitHub, bấm Đồng bộ GitHub .tex

% Bài 2. Tập hợp và các phép toán trên tập hợp


\begin{ex}
 Cho hai tập hợp $A=\{{-6, 3, -4, -1}\}$ và $B=\{{1, 2, -10, -7, -2}\}$. Tìm tập hợp $A\cap B$. \\ 
\choice
{ $\{{-6; 3; -4; -1}\}$ }
   { \True $\emptyset$ }
     { $\{{1; 2; 3; -10; -2; -7; -6; -4; -1}\}$ }
    { $\{{1; 2; -10; -7; -2}\}$ }
\loigiai{ 
  
 }\end{ex}

\begin{ex}
 Cho hai tập hợp $A=\{{1, 3, 5, -5, -4}\}$ và $B=\{{2, 3, 7, -8, -3, -1}\}$. Tìm tập hợp $A\cup B$. \\ 
\choice
{ $\{{2; 7; -8; -3; -1}\}$ }
   { $\{{1; -5; -4; 5}\}$ }
     { $\{{3}\}$ }
    { \True $\{{1; 2; 3; 5; 7; -8; -5; -4; -3; -1}\}$ }
\loigiai{ 
  
 }\end{ex}

\begin{ex}
 Cho hai tập hợp $M=\{x\in \mathbb{Z}|-1 \le x < 5\}$ và $N=\{x\in \mathbb{Z}|-4 < x \le 2\}$. Tìm tập hợp $M \backslash N$.\\ 
\choice
{ $\{{-3; -2}\}$ }
   { $\{{0; 1; 2; 3; 4; -2; -3; -1}\}$ }
     { $\{{0; 1; 2; -1}\}$ }
    { \True $\{{3; 4}\}$ }
\loigiai{ 
 Ta có: $M=\{{0, 1, 2, 3, 4, -1}\}$ và $N=\{{0, 1, 2, -1, -3, -2}\}$. Do đó $M \backslash N =\{{3, 4}\}$. 
 }\end{ex}

\begin{ex}
 Cho hai tập hợp ${A}$ và ${B}$ biết $n(A)=27,n(B)=29, n(A\cup B)=41$. Tính $n(A\cap B)$. 
\choice
{ ${37}$ }
   { ${14}$ }
     { ${12}$ }
    { \True ${15}$ }
\loigiai{ 
 $n(A\cap B)=n(A)+n(B)-n(A\cup B) =27+29-41=15$. 
 }\end{ex}

\begin{ex}
 Cho hai tập hợp ${C}$ và ${M}$ biết $n(C)=23,n(M)=25, n(C\cup M)=34$. Tính $n(C\cap M)$. 
\choice
{ ${38}$ }
   { ${48}$ }
     { \True ${14}$ }
    { ${2}$ }
\loigiai{ 
 $n(C\cap M)=n(C)+n(M)-n(C\cup M) =23+25-34=14$. 
 }\end{ex}

\begin{ex}
 Cho hai tập hợp ${N}$ và ${D}$ biết $n(N)=22,n(D)=21, n(N\cup D)=33$. Tính $n(N\cap D)$. 
\choice
{ ${11}$ }
   { ${12}$ }
     { \True ${10}$ }
    { ${1}$ }
\loigiai{ 
 $n(N\cap D)=n(N)+n(D)-n(N\cup D) =22+21-33=10$. 
 }\end{ex}

\begin{ex}
 Cho hai tập hợp ${E}$ và ${M}$ biết $n(E)=24,n(M)=32, n(E\cup M)=41$. Tính $n(E\cap M)$. 
\choice
{ \True ${15}$ }
   { ${9}$ }
     { ${45}$ }
    { ${56}$ }
\loigiai{ 
 $n(E\cap M)=n(E)+n(M)-n(E\cup M) =24+32-41=15$. 
 }\end{ex}

\begin{ex}
 Cho hai tập hợp ${C}$ và ${N}$ biết $n(C)=23,n(N)=30, n(C\cup N)=40$. Tính $n(C\cap N)$. 
\choice
{ ${53}$ }
   { \True ${13}$ }
     { ${50}$ }
    { ${10}$ }
\loigiai{ 
 $n(C\cap N)=n(C)+n(N)-n(C\cup N) =23+30-40=13$. 
 }\end{ex}

\begin{ex}
 Cho hai tập hợp ${N}$ và ${E}$ biết $n(N)=17,n(E)=21, n(N\cup E)=29$. Tính $n(N\cap E)$. 
\choice
{ ${12}$ }
   { ${34}$ }
     { \True ${9}$ }
    { ${4}$ }
\loigiai{ 
 $n(N\cap E)=n(N)+n(E)-n(N\cup E) =17+21-29=9$. 
 }\end{ex}

\begin{ex}
 Cho hai tập hợp ${E}$ và ${M}$ biết $n(E)=28,n(M)=32, n(E\cup M)=46$. Tính $n(E\cap M)$. 
\choice
{ ${60}$ }
   { ${18}$ }
     { \True ${14}$ }
    { ${48}$ }
\loigiai{ 
 $n(E\cap M)=n(E)+n(M)-n(E\cup M) =28+32-46=14$. 
 }\end{ex}

\begin{ex}
 Cho hai tập hợp ${A}$ và ${B}$ biết $n(A)=19,n(B)=18, n(A\cup B)=30$. Tính $n(A\cap B)$. 
\choice
{ ${37}$ }
   { ${1}$ }
     { ${11}$ }
    { \True ${7}$ }
\loigiai{ 
 $n(A\cap B)=n(A)+n(B)-n(A\cup B) =19+18-30=7$. 
 }\end{ex}

\begin{ex}
 Cho hai tập hợp ${A}$ và ${C}$ biết $n(A)=21,n(C)=33, n(A\cup C)=41$. Tính $n(A\cap C)$. 
\choice
{ ${8}$ }
   { ${48}$ }
     { \True ${13}$ }
    { ${54}$ }
\loigiai{ 
 $n(A\cap C)=n(A)+n(C)-n(A\cup C) =21+33-41=13$. 
 }\end{ex}

\begin{ex}
 Cho hai tập hợp ${N}$ và ${A}$ biết $n(N)=22,n(A)=23, n(N\cup A)=38$. Tính $n(N\cap A)$. 
\choice
{ ${16}$ }
   { ${15}$ }
     { ${45}$ }
    { \True ${7}$ }
\loigiai{ 
 $n(N\cap A)=n(N)+n(A)-n(N\cup A) =22+23-38=7$. 
 }\end{ex}
\begin{ex}
 Biết rằng lớp 12A1 có ${20}$ bạn biết chơi bóng đá và ${18}$ bạn biết chơi bóng chuyền. Trong số các bạn biết chơi bóng đá hoặc biết chơi bóng chuyền có ${6}$ bạn biết chơi cả hai. Lớp vẫn còn ${4}$ bạn không biết chơi bóng đá và không biết chơi bóng chuyền. Xét tính đúng-sai của các khẳng định sau.
\choiceTFt
{ \True Số học sinh chỉ biết chơi bóng đá là ${14}$ }
   { \True Số học sinh chỉ biết chơi bóng chuyền là ${12}$ }
     { \True Số học sinh biết chơi bóng đá hoặc biết chơi bóng chuyền là ${32}$ }
    { Tổng số học sinh của lớp 12A1 là ${39}$ }
\loigiai{ 
 

 a) Khẳng định đã cho là khẳng định đúng.

 Số học sinh chỉ biết chơi bóng đá là: $20-6=14$.

b) Khẳng định đã cho là khẳng định đúng.

 Số học sinh chỉ biết chơi bóng chuyền là: $18-6=12$.

c) Khẳng định đã cho là khẳng định đúng.

 Số học sinh biết chơi bóng đá hoặc biết chơi bóng chuyền là: $20+18-6=32$.

d) Khẳng định đã cho là khẳng định sai.

 Số học sinh biết chơi bóng đá hoặc biết chơi bóng chuyền là: $20+18-6=32$.

Tổng số học sinh của lớp là: $32+4=36$.

 
 }\end{ex}

\begin{ex}
 Biết rằng lớp 11B6 có ${20}$ bạn biết chơi cờ vua và ${22}$ bạn biết chơi cờ tướng. Trong số các bạn biết chơi cờ vua hoặc biết chơi cờ tướng có ${6}$ bạn biết chơi cả hai. Lớp vẫn còn ${5}$ bạn không biết chơi cờ vua và không biết chơi cờ tướng. Xét tính đúng-sai của các khẳng định sau.
\choiceTFt
{ \True Số học sinh chỉ biết chơi cờ vua là ${14}$ }
   { Số học sinh chỉ biết chơi cờ tướng là ${17}$  }
     { \True Số học sinh biết chơi cờ vua hoặc biết chơi cờ tướng là ${36}$ }
    { Tổng số học sinh của lớp 11B6 là ${43}$ }
\loigiai{ 
 

 a) Khẳng định đã cho là khẳng định đúng.

 Số học sinh chỉ biết chơi cờ vua là: $20-6=14$.

b) Khẳng định đã cho là khẳng định sai.

 Số học sinh chỉ biết chơi cờ tướng là: $22-6=16$.

c) Khẳng định đã cho là khẳng định đúng.

 Số học sinh biết chơi cờ vua hoặc biết chơi cờ tướng là: $20+22-6=36$.

d) Khẳng định đã cho là khẳng định sai.

 Số học sinh biết chơi cờ vua hoặc biết chơi cờ tướng là: $20+22-6=36$.

Tổng số học sinh của lớp là: $36+5=41$.

 
 }\end{ex}

\begin{ex}
 Biết rằng lớp 10A4 có ${17}$ bạn thích nghe nhạc và ${20}$ bạn thích xem phim. Trong số các bạn thích nghe nhạc hoặc thích xem phim có ${7}$ bạn thích cả hai. Lớp vẫn còn ${7}$ bạn không thích nghe nhạc và không thích xem phim. Xét tính đúng-sai của các khẳng định sau.
\choiceTFt
{ \True Số học sinh chỉ thích nghe nhạc là ${10}$ }
   { Số học sinh chỉ thích xem phim là ${16}$  }
     { \True Số học sinh thích nghe nhạc hoặc thích xem phim là ${30}$ }
    { Tổng số học sinh của lớp 10A4 là ${39}$ }
\loigiai{ 
 

 a) Khẳng định đã cho là khẳng định đúng.

 Số học sinh chỉ thích nghe nhạc là: $17-7=10$.

b) Khẳng định đã cho là khẳng định sai.

 Số học sinh chỉ thích xem phim là: $20-7=13$.

c) Khẳng định đã cho là khẳng định đúng.

 Số học sinh thích nghe nhạc hoặc thích xem phim là: $17+20-7=30$.

d) Khẳng định đã cho là khẳng định sai.

 Số học sinh thích nghe nhạc hoặc thích xem phim là: $17+20-7=30$.

Tổng số học sinh của lớp là: $30+7=37$.

 
 }\end{ex}

\begin{ex}
 Biết rằng lớp 10B9 có ${20}$ bạn biết chơi bóng đá và ${22}$ bạn biết chơi bóng chuyền. Trong số các bạn biết chơi bóng đá hoặc biết chơi bóng chuyền có ${8}$ bạn biết chơi cả hai. Lớp vẫn còn ${5}$ bạn không biết chơi bóng đá và không biết chơi bóng chuyền. Xét tính đúng-sai của các khẳng định sau.
\choiceTFt
{ \True Số học sinh chỉ biết chơi bóng đá là ${12}$ }
   { Số học sinh chỉ biết chơi bóng chuyền là ${17}$  }
     { \True Số học sinh biết chơi bóng đá hoặc biết chơi bóng chuyền là ${34}$ }
    { \True Tổng số học sinh của lớp 10B9 là ${39}$ }
\loigiai{ 
 

 a) Khẳng định đã cho là khẳng định đúng.

 Số học sinh chỉ biết chơi bóng đá là: $20-8=12$.

b) Khẳng định đã cho là khẳng định sai.

 Số học sinh chỉ biết chơi bóng chuyền là: $22-8=14$.

c) Khẳng định đã cho là khẳng định đúng.

 Số học sinh biết chơi bóng đá hoặc biết chơi bóng chuyền là: $20+22-8=34$.

d) Khẳng định đã cho là khẳng định đúng.

 Số học sinh biết chơi bóng đá hoặc biết chơi bóng chuyền là: $20+22-8=34$.

Tổng số học sinh của lớp là: $34+5=39$.

 
 }\end{ex}

\begin{ex}
 Biết rằng lớp 12B1 có ${18}$ bạn thích nghe nhạc và ${23}$ bạn thích xem phim. Trong số các bạn thích nghe nhạc hoặc thích xem phim có ${8}$ bạn thích cả hai. Lớp vẫn còn ${7}$ bạn không thích nghe nhạc và không thích xem phim. Xét tính đúng-sai của các khẳng định sau.
\choiceTFt
{ \True Số học sinh chỉ thích nghe nhạc là ${10}$ }
   { \True Số học sinh chỉ thích xem phim là ${15}$ }
     { \True Số học sinh thích nghe nhạc hoặc thích xem phim là ${33}$ }
    { Tổng số học sinh của lớp 12B1 là ${42}$ }
\loigiai{ 
 

 a) Khẳng định đã cho là khẳng định đúng.

 Số học sinh chỉ thích nghe nhạc là: $18-8=10$.

b) Khẳng định đã cho là khẳng định đúng.

 Số học sinh chỉ thích xem phim là: $23-8=15$.

c) Khẳng định đã cho là khẳng định đúng.

 Số học sinh thích nghe nhạc hoặc thích xem phim là: $18+23-8=33$.

d) Khẳng định đã cho là khẳng định sai.

 Số học sinh thích nghe nhạc hoặc thích xem phim là: $18+23-8=33$.

Tổng số học sinh của lớp là: $33+7=40$.

 
 }\end{ex}
\begin{ex}
 Lớp 12A3 có tổng cộng ${36}$ học sinh, các học sinh này đều thích xem phim rạp hoặc thích xem ca nhạc. Có ${20}$ học sinh thích xem phim rạp (trong số này có các học sinh thích xem ca nhạc) và ${24}$ học sinh thích xem ca nhạc (trong số này có các học sinh thích xem phim rạp). Hỏi lớp 12A3 có bao nhiêu học sinh thích xem phim rạp và thích xem ca nhạc.
\shortans[4]{${8}$}

\loigiai{ 
 Số học sinh chỉ thích xem phim rạp và không thích xem ca nhạc là: $36-24=12$.

Số học sinh thích xem phim rạp và thích xem ca nhạc là: $20-12=8$. 
 }\end{ex}

\begin{ex}
 Mỗi học sinh của lớp 10A5 đều thích nấu ăn hoặc thích đọc sách. Biết rằng lớp có ${19}$ bạn thích nấu ăn (trong số này có các bạn thích đọc sách), có ${17}$ bạn thích đọc sách (trong số này có các bạn thích nấu ăn) và có ${5}$ bạn thích nấu ăn và thích đọc sách. Hỏi lớp 10A5 có tổng cộng bao nhiêu học sinh?
\shortans[4]{${31}$}

\loigiai{ 
 Gọi ${A}$ là tập hợp các bạn thích nấu ăn, ta có: $n(A)=19$.

Gọi ${B}$ là tập hợp các bạn thích đọc sách, ta có: $n(B)=17$.

$A\cap B$ là tập hợp các bạn thích nấu ăn và thích đọc sách, ta có: $n(A\cap B)=5$.

$C=A\cup B$ là tập hợp tất cả các bạn của lớp 10A5, ta có: $n(C)=19+17-5=31$. 
 }\end{ex}

\begin{ex}
 Biết rằng lớp 10A6 có ${17}$ bạn thích nấu ăn và ${21}$ bạn thích đọc sách. Trong số các bạn thích nấu ăn hoặc thích đọc sách có ${5}$ bạn thích cả hai. Lớp vẫn còn ${5}$ bạn không thích nấu ăn và không thích đọc sách. Hỏi lớp 10A6 có tổng cộng bao nhiêu học sinh?
\shortans[4]{${38}$}

\loigiai{ 
 Gọi ${A}$ là tập hợp các bạn thích nấu ăn, ta có: $n(A)=17$.

Gọi ${B}$ là tập hợp các bạn thích đọc sách, ta có: $n(B)=21$.

$A\cap B$ là tập hợp các bạn thích nấu ăn và thích đọc sách, ta có: $n(A\cap B)=5$.

$A\cup B$ là tập hợp tất cả các bạn thích nấu ăn hoặc thích đọc sách, ta có: $n(A\cup B)=17+21-5=33$.

Tổng số học sinh của lớp là: $33+5=38$. 
 }\end{ex}

\begin{ex}
 Lớp 12B6 có ${17}$ bạn biết chơi bóng đá, ${15}$ bạn biết chơi cờ tướng, ${14}$ bạn biết chơi cờ vua, ${4}$ bạn biết chơi bóng đá và biết chơi cờ tướng, ${4}$ bạn biết chơi cờ tướng và biết chơi cờ vua, ${6}$ bạn biết chơi bóng đá và biết chơi cờ vua và ${3}$ biết chơi cả ba môn. Hỏi lớp 12B6 có tất cả bao nhiêu học sinh biết ít nhất một môn thể thao?
\shortans[4]{35}

\loigiai{ 
 Gọi A là tập hợp các bạn biết chơi bóng đá, B là tập hợp các bạn biết chơi cờ tướng, C là tập hợp các bạn biết chơi cờ vua.

Ta có: $n(A)=17, n(B)=15, n(C)=14$.

$n(A\cap B)=4, n(A\cap C)=6, n(B\cap C)=4$.

$n(A\cap B \cap C)=3$.

Số học sinh biết ít nhất một môn thể thao là:

$n(A\cup B \cup C)=n(A)+n(B)+n(C)+n(A\cap B \cap C)-n(A\cap B)-n(A\cap C)-n(B\cap C)$

$=17+15+14+3-4-6-4=35$.

Đáp án: 35 
 }\end{ex}

\begin{ex}
 Lớp 12A9 có ${20}$ bạn thích đọc sách, ${21}$ bạn thích đi du lịch, ${20}$ bạn thích nấu ăn, ${6}$ bạn thích đọc sách và thích đi du lịch, ${7}$ bạn thích đi du lịch và thích nấu ăn, ${7}$ bạn thích đọc sách và thích nấu ăn, ${4}$ có cả ba sở thích. Hỏi lớp 12A9 có tất cả bao nhiêu học sinh chỉ có đúng một sở thích trong các sở thích trên?
\shortans[4]{33}

\loigiai{ 
 Gọi A là tập hợp các bạn thích đọc sách, B là tập hợp các bạn thích đi du lịch, C là tập hợp các bạn thích nấu ăn.

Ta có: $n(A)=20, n(B)=21, n(C)=20$.

$n(A\cap B)=6, n(A\cap C)=7, n(B\cap C)=7$.

$n(A\cap B \cap C)=4$.

Số học sinh chỉ thích đọc sách là: $20-6-7+4=11$

Số học sinh chỉ thích đi du lịch là: $21-6-7+4=12$

Số học sinh chỉ thích nấu ăn là: $20-7-7+4=10$

Số học sinh chỉ có đúng một sở thích: $11+12+10=33$.

Đáp án: 33 
 }\end{ex}
